\section{Introduction}

\subsection{Motivation}

The proliferation of large language models has made text paraphrasing trivially accessible. A student submitting coursework, a content farm scaling production, or a disinformation operator laundering provenance can each iteratively paraphrase a document through multiple LLMs, producing text that retains the original meaning while evading plagiarism detectors and authorship attribution systems. This creates urgent challenges for academic integrity, forensic linguistics, and AI transparency: if an LLM can rewrite text until no trace of the original author remains, traditional authorship verification collapses.

The threat is not hypothetical. Deployed detectors such as Turnitin report that a substantial fraction of student writing now contains AI-generated material at scale \cite{turnitin2024aiwriting}, and independent studies show that neural paraphrasers can evade state-of-the-art AI text detectors \cite{krishna2023dipper}, with iterative rewriting compounding the evasion effect. Emerging regulation such as the EU AI Act \cite{euaiact2024} expects AI-generated content to be labelled or otherwise flagged; those obligations are undermined whenever provenance signals can be stripped by rewriting. The consequences touch academic integrity offices deciding whether to escalate a suspected submission, newsrooms attempting to verify the origin of a leaked quote, and platforms enforcing synthetic-media policies. In each of these settings, a defender cannot simply ask \emph{was this text machine-generated?}: they must ask \emph{how much of the original author is still present in this version?} Yet most detection research still treats paraphrasing as a single-step perturbation, measuring whether a classifier can survive one rewrite. The more fundamental question---how authorial identity \emph{decays} across multiple generations of paraphrasing, which specific linguistic features erode first, and how quickly a usable forensic signal disappears---remains largely unexplored.

\subsection{Conceptual Framing}

We frame this problem through the Ship of Theseus paradox: if you replace every plank of a ship one by one, is the result still the same ship? Analogously, if an LLM replaces every word, restructures every sentence, and sheds every proper noun while preserving the original meaning, is the result still the original author's work? We map the philosophical framework onto measurable NLP signals: \emph{lexical planks} (individual words, measured by BLEU), \emph{structural hull} (grammatical patterns, measured by POS tag distributions), \emph{entity cargo} (factual anchors, measured by NER), and \emph{semantic meaning} (measured by BERTScore). This mapping allows us to quantify which ``planks'' are replaced first and at what rate.

\subsection{Research Gap}

Prior work on authorship attribution under adversarial conditions \cite{koppel2011authorship} demonstrates that paraphrasing degrades attribution accuracy, but does not quantify the \emph{rate} or \emph{order} of feature decay across iterations. Work on AI text detection \cite{mitchell2023detectgpt} focuses on binary classification (human vs.\ machine) rather than measuring how iterative rewriting erodes specific linguistic features at different rates. Neural paraphrase typologies \cite{wahle2023paraphrase} characterize the \emph{types} of changes paraphrasers make, but do not track how those changes accumulate across multiple generations.

The distinction between single-step and iterative paraphrasing is fundamental. Single-step studies ask a threshold question: ``can a detector survive one attack?'' Iterative studies ask a phase-transition question: ``at what depth does the authorial signal become unrecoverable?'' Traditional authorship attribution methods---from Mosteller and Wallace's Bayesian analysis of \emph{The Federalist Papers} to Burrows' Delta and its successors---were designed for static corpora where text is fixed. They provide no framework for measuring how stylometric features degrade under systematic, iterative transformation. Our work bridges this gap by treating authorial identity not as a binary property (present or absent) but as a continuously decaying signal whose components erode at measurably different rates.

Our work fills this gap by providing a multi-signal, multi-iteration forensic analysis that measures \emph{which} features decay, \emph{how fast}, and \emph{whether different paraphrasers leave identifiable signatures}. Treating identity as a decaying signal rather than a binary label lets us describe the forensic landscape at each iteration depth, and lets practitioners reason explicitly about which detectors remain trustworthy after one, two, or three rewrites.

\subsection{Research Questions}

We investigate three questions using the Ship of Theseus Paraphrased Corpus \cite{tripto2024ship}, each targeting a different aspect of the forensic problem: RQ1 characterizes \emph{what} decays and how fast, RQ2 identifies the \emph{threshold} at which identity is lost, and RQ3 asks whether the transformation tool itself can be identified---akin to determining which shipyard replaced the planks.

\begin{itemize}
    \item \textbf{RQ1:} Which linguistic features decay first---style (POS distributions, punctuation) or content (named entities, factual anchors)?
    \item \textbf{RQ2:} At what paraphrase iteration does authorship attribution effectively fail?
    \item \textbf{RQ3:} Do different paraphrasing models leave identifiable forensic fingerprints?
\end{itemize}

\textbf{Contributions.} We present a multi-signal forensic analysis combining lexical metrics (BLEU, ROUGE-L), semantic metrics (BERTScore), stylometric features (POS tag distributions), and content features (NER entity preservation) across seven paraphraser variants, seven diverse datasets, and three paraphrase iterations. We find that (1) different feature types exhibit \emph{layered erosion}, with lexical surface vanishing first, entities next, and grammatical structure last; (2) content features (named entities) are 5.8$\times$ more sensitive to paraphrasing intensity than stylistic features (POS distributions); (3) entities are predominantly \emph{dropped} rather than replaced by novel entities; and (4) each paraphraser occupies a distinct position in the style-vs-content preservation space, with paraphraser configuration mattering more than model architecture.
